\section{Boundary geometry and the Milnor representative}\label{sec:boundary-geometry}

The tube calculation must describe a genuine local Milnor page.
We choose a larger representative whose inner and exterior
parts admit the required continuations. This section proves its
boundary regularity, constructs compatible changes of model on the
inner and exterior pieces, and proves isolation. The next section
uses those models to prove acyclicity and then compares this
representative with an ordinary small Milnor tube.

Fix \(M>3\), with \(M^{a+b}>1+k\max(a,b)\), and for
\(t>0\) set
\[
Y_t=\{|y|\le Mt^\alpha,\ |z|\le Mt^\beta,\ |w|\le t\}.
\]
Choose a sufficiently small base ball \(\mathcal B\) for both
plane tubes. Then choose \(\rho>0\), \(\varepsilon>0\),
\(\delta>0\), \(\lambda>0\), and \(s>0\), in that order, with
\begin{equation}
0<\delta<\varepsilon/4,\qquad 0<\lambda\ll\rho,\qquad
\sup_{Y_\lambda}|P|<\delta/4,
\label{eq:tube-3}
\end{equation}
and with every root of \(R_s(p,\cdot)\), for \(p\in\mathcal B\),
in \(|w|<\lambda/4\). The last condition follows because the roots
are \(s\) times the roots of \(R(p,\cdot)\) on the compact base cut.
Our larger closed total tube is
\begin{equation}
\widehat\T_s=
\{p\in\mathcal B,\ |h|\le\varepsilon,\ (y,z,w)\in Y_\rho,\
yz=R_s(p,w),\ |h+P|\le\delta\},
\label{eq:tube-4}
\end{equation}
and \(\T_s=f_s^{-1}(c)\cap\widehat\T_s\) for \(|c|=\delta\).
Write \(X_s\) for the same cut set with the final inequality
\(|h+P|\le\delta\) omitted. Its outer lateral strata are formed by
the base-ball, \(|h|=\varepsilon\), and \(\partial Y_\rho\) faces
and their active corners.
Let \(W\) be the part of \(\T_s\) in \(Y_\lambda\), let
\(M_{\rm ext}\) be the part of \(\T_s\) outside the interior of
\(Y_\lambda\), and let \(J\) be their
bicollared interface. Hats denote the corresponding total disk pieces.
The auxiliary interface \(J\) is an interior cut of \(X_s\), rather
than one of its outer lateral strata.

\subsection{Regularity on the outer boundary}

\begin{lemma}[Outer boundary regularity]\label{lem:boundary}
For every sufficiently small fixed \(s>0\), the map \(f_s\) on
the compact outer lateral boundary of \(X_s\) over
\(\{|v|\le\delta\}\) is proper, and its restriction to every
stratum has real rank two throughout that disk.
In particular, all active lateral faces and corners of \eqref{eq:tube-4}
are independent of the horizontal face over \(|f_s|=\delta\).
The horizontal boundary map is proper on the union of lateral
strata and a submersion on each stratum.
\end{lemma}
\begin{proof}
We first estimate the vertical faces and corners, then prove
uniform regularity on the remaining outer strata.

At \(s=0\) the vertical equation is \(yz=w^k\). Its link on
\(\partial Y_t\) is the cyclic quotient of a 3-sphere, hence a lens
space. The cone is smooth away from its vertex. The \(y\)- and
\(z\)-faces cannot meet: if both were active, then
\(|yz|=M^2t^k>|w|^k\). On the \(y\)-face,
\begin{equation}
\left|\frac{z^b}{y^a}\right|
\le M^{-(a+b)}\left(\frac{|w|}{t}\right)^{kb}<1.
\label{eq:tube-5}
\end{equation}
After eliminating \(z=w^k/y\), the free derivative is
\(P_w=y^a+(kb+1)z^b\), which is nonzero by the stronger choice
of \(M\). The \(z\)-face is symmetric, with
\(P_w=z^b+(ka+1)y^a\). These estimates persist for small
positive \(s\) by uniform convergence of \(R_s\) and
\(\partial_wR_s\) on the fixed outer face.
At a \(y/w\) corner the bound
\begin{equation}
 |P|\ge t^{1+a\alpha}(M^a-M^{-b})>0
\label{eq:tube-6}
\end{equation}
holds; similarly at a \(z/w\) corner. Choose
\(\varepsilon+\delta\) below both lower bounds at \(t=\rho\).
On the \(w\)-face, near a zero of \(P\) with \(w\ne0\),
use \(y\) as the free coordinate and \(z=w^k/y\).
Then \(dP(\partial_{\log y})=w(ay^a-bz^b)
=w(a+b)y^a\ne0\) at \(P=0\).
On the compact outer \(w\)-face, this nonvanishing persists in
one uniform small-value neighborhood; choose
\(\varepsilon+\delta\) within it.
The quotient expression \(\widetilde P=UV(U^{ka}+V^{kb})\)
also shows submersivity of \(P\) in the nonvertex cone interior:
if \(U,V\ne0\), its two partials cannot vanish simultaneously
because that would imply \((ka+1)(kb+1)=1\);
on each axis the transverse partial is nonzero.
These calculations cover the vertical faces and corners near the
small values used in \eqref{eq:tube-4}.

To finish the regularity claim on the genuine outer strata, consider
a sequence approaching the cone vertex while \(s\to0\), and let
\(\varrho=\max(|y|^{1/\alpha},|z|^{1/\beta},|w|)\).
If \(\varrho/s\to\infty\), weighted rescaling by \(\varrho\) gives smooth
convergence at a nonvertex cone point; all lower Weierstrass terms
vanish because \(A_j\) is bounded and \((s/\varrho)^{k-j}\to0\).
If \(\varrho/s\) is bounded, rescale by \(s\), with
\((Y,Z,u)=(y/s^\alpha,z/s^\beta,w/s)\); the limit is the smooth
model \(YZ=R(p,u)\). For \((Y,Z)\ne(0,0)\)
the base projection is a submersion; at \(Y=Z=0\) its differential
contains the tangent image of the smooth root graph.
The specified \(h,d\) Milnor cuts have the needed transversality
on \(\partial\mathcal B\), including corners with \(|h|=\varepsilon\).
For the actual function \(f_s\), these arguments give the stated
outer lateral regularity over the whole value disk. Indeed
\(\lvert P\rvert\le\varepsilon+\delta\) whenever
\(\lvert f_s\rvert\le\delta\), so all the outer vertical-face
estimates apply throughout that disk, also at their intersections
with base faces because the free vertical directions keep the base
fixed. On an outer stratum with no vertical face, a hypothetical
failure with unbounded \(\varrho/s\) would contradict submersivity of
\(P\) on the nonvertex cone after rescaling. If \(\varrho/s\) is bounded,
then \(P=s^{1+a\alpha}u(Y^a+Z^b)=o(1)\) in the smooth rescaled
model. The face \(\lvert h\rvert=\varepsilon\) is then absent over
\(\lvert f_s\rvert\le\delta<\varepsilon/4\). On the remaining
base-ball face, the boundary fibrations of \(h\) and \(d\) over
their full small value disks give rank two, respectively where
\((Y,Z)\ne(0,0)\) and on the smooth root graph. Compactness and
the preceding smooth convergence preserve these ranks for small
\(s\). This assertion concerns the outer cuts of \(X_s\); it
does not assert disk-wide submersivity on the auxiliary \(J\).
The remaining interior critical-point possibility is excluded in
Lemma~\ref{lem:isolation}. All statements here are open in compact
families of smooth strata, so one uniform upper bound on \(s\)
works for all strata.

\end{proof}

\subsection{Compatible continuations and the lens circle}

The inner and exterior computations use different models. To glue
classes computed in them, we must carry the same oriented lens circle
over each labeled base loop. This retains the labeling of the
interface tori; in Section~\ref{sec:integral-lattice} it will identify
the full finite-map transfer.

For \(|c|=\delta\), write \(\Hh_c=\mathcal B\cap\{h=c\}\).
Choose a smooth \(h\)-trivialization of pairs, respecting the
base-ball boundary, \(\Phi(p_0,q)\), over a disk about \(c\)
containing \(\{|q-c|\le\delta/4\}\), with
\(p_0\in\Hh_c\), \(\Phi(p_0,c)=p_0\), and
\(h(\Phi(p_0,q))=q\). Set \(w_0=\lambda/2\).

\begin{lemma}[Compatible continuations]\label{lem:continuations}
For sufficiently small \(s>0\), the exterior total/page pairs continue
to \(s=0\). On the inner part, the sublevel totals and
horizontal-boundary fibrations continue from \(h+P\) to \(h\)
through \(h+\theta P\), \(0\le\theta\le1\). These continuations preserve the complex
value on the horizontal boundary and agree on \(J\) up to isotopy
through their two-parameter rectangle.

They can be chosen to preserve the label \(p_0\), the coordinate
\(w_0\), and the exact positive \(y\)-angle \(\phi\) on the
interface circles characterized by
\[
 p=\Phi(p_0,q),\qquad w=w_0,\qquad
 y=M\lambda^\alpha e^{i\phi},\qquad h(p)+\theta P(y,z,w)=c.
\]
Here \(z=R_s(p,w_0)/y\), and \(q\) is the unique value near \(c\)
that satisfies these equations. The inner continuation is an isotopy
of sublevel totals; value preservation is required on its horizontal
boundary, not through its critical region.
\end{lemma}
\begin{proof}
At the inner interface, \eqref{eq:tube-3} and \(|h+\theta P|=\delta\) imply
\begin{equation}
3\delta/4\le |h|\le5\delta/4<\varepsilon
\quad(0\le\theta\le1).
\label{eq:tube-7}
\end{equation}
Thus \(h\) is regular there, including on its base-ball cut, and
the \(|h|=\varepsilon\) cut is absent. On the inner \(w\)-face,
\eqref{eq:tube-7} makes \(h\) regular, and \(\lambda\) can be
decreased until \(P\) is small in \(C^1\) there.
The zero-parameter cone is
smooth along this interface, so the same assertion holds throughout
the \((s,\theta)\) rectangle for small \(s\). We use regularity
of the horizontal boundary to isotope the \emph{sublevel totals}
as manifolds with corners; we do not assert that \(h\) is submersive
through every value in the disk at \(\theta=0\).

We first construct the prescribed circle family. For fixed
\((p_0,\phi,\theta,s)\), solve
\[
q+\theta P\left(y,\frac{R_s(\Phi(p_0,q),w_0)}y,w_0\right)=c
\]
for \(q\) near \(c\), then set
\(p=\Phi(p_0,q)\) and \(z=R_s(p,w_0)/y\).
The root bound gives
\(\lvert R_s(p,w_0)\rvert\le(3\lambda/4)^k\), so these circles
lie on the \(y\)-face and strictly away from the \(z\)- and
\(w\)-faces. Thus the small bound for \(P\) applies there.
The solution is unique and smooth by the real implicit-function
theorem: \(P\) is uniformly small there, and its real derivative
\(D_qP\to0\) as \(s\to0\), since every
\(p\)-dependent term of \(R_s(p,w_0)\) contains a positive
power of \(s\). Decrease \(s\) until \(\lVert D_qP\rVert<1/2\)
uniformly on this neighborhood. This lift keeps the \(h\)-page label \(p_0\),
\(w_0\), and the exact oriented \(y\)-angle \(\phi\) fixed;
the physical \(p\) may move. At \(\theta=0\), the \(s\)-motion
keeps \(p_0,y,w_0\) fixed and only changes
\(z=R_s(p_0,w_0)/y\). Hence the family of anchored
circles is \(\Hh_c\times S^1\), with no base-to-circle shear.
To make this family the restriction of the boundary isotopy, use
the same coordinates near its compact trace and put
\[
 \mathcal P_s(p_0,q,y,w)
 =P\left(y,\frac{R_s(\Phi(p_0,q),w)}{y},w\right).
\]
The real operator \(I+\theta D_q\mathcal P_s\) is uniformly
invertible there. The local lifts that fix \(p_0,y,w\) have
\[
 \dot q_\theta=-(I+\theta D_q\mathcal P_s)^{-1}\mathcal P_s,
 \qquad
 \dot q_s=-(I+\theta D_q\mathcal P_s)^{-1}
                  \theta\,\partial_s\mathcal P_s.
\]
Differentiating \(q+\theta\mathcal P_s\) verifies that both lifts
preserve its complex value. They are tangent to the \(y\)-face
and to the base-ball face, because \(y\) and the pair-trivialization
label \(p_0\) are fixed. Their restrictions to the circles are
exactly the derivatives of the displayed implicit solution.

We extend these local lifts to the sublevel continuations.
Near \(J\), vary \(\lambda\) slightly to obtain bicollars.
For the \(\theta\)-motion the ambient \(C_s\) is fixed. On a collar
of \(|h+\theta P|=\delta\), the map consisting of \(h+\theta P\)
and the active lateral defining functions is a submersion. Choose
a parameter field \(V\) tangent to every active lateral face and
satisfying \(d(h+\theta P)(V)+P=0\) there. Near the compact anchor
trace, patch it with the local \(\theta\)-lift above by a cutoff
equal to one on that trace. The lift equation and all active-face
tangency conditions are affine, so this patch remains a valid lift.
Extend only its spatial correction by a cutoff supported in that
collar and zero near the critical region, retaining the parameter
direction, and integrate.
This isotopes the smooth sublevel totals and preserves the complex
value on their horizontal boundary. For the \(s\)-motion on the exterior,
identify the smooth ambient family away from the cone vertex and make
the same horizontal-collar correction, patched with the displayed
local \(s\)-lift near the anchors. Uniqueness of the flows makes
their restrictions exactly the anchored family, fixing
\(p_0,w_0,\phi\). On the interface rectangle
\eqref{eq:tube-7} keeps \(h\) regular, so its collar lifts exist along
every edge and path homotopy gives isotopic boundary identifications
relative to these anchors. Locally this is explicit in the coordinates
\((p_0,q+\theta\mathcal P_s,y,w,\theta,s)\), in which the two
anchored lifts are parameter coordinate fields.
At \(s=0\), the cone is independent of \(p\), so on the entire
page interface this normalization can be chosen explicitly as
\[
 (p_0,v)\longmapsto\bigl(\Phi(p_0,c-\theta P(v)),v\bigr),
 \qquad v=(y,z,w)\in L=\{yz=w^k\}\cap\partial Y_\lambda.
\]
The bound \(|P(v)|<\delta/4\) keeps this formula in the chosen
trivialization. It agrees with the anchored circles and identifies
the interface with \(\Hh_c\times L\) while retaining the same
base label. No value-preserving lift is used through the critical
inner total.
\end{proof}

\subsection{Isolation in the chosen tube}

\begin{lemma}[Uniform critical-point isolation]\label{lem:isolation}
For sufficiently large \(a,b\) as in Theorem~\ref{thm:tube},
the only critical point of \(f_s\) in \eqref{eq:tube-4} for all sufficiently
small \(s>0\) is the origin.
\end{lemma}
\begin{proof}
The gradient ideal of the isolated plane singularity \(h\) contains
a power of the maximal ideal. For this \(h\), use the same
\(L_0=4\) supplied by \eqref{eq:plane-gradient}: on each nonconstant
analytic arc \(p\to0\),
\(\nu(h_p)\le L_0\nu(p)\).
For fixed \(s\), at a constrained critical point in the ambient
coordinates, write the spatial differential equation
\[
 d(h+P)=\chi_{\rm orig}\,d(yz-R_s).
\]
After the change \(y=s^\alpha Y\), \(z=s^\beta Z\), \(w=su\),
put \(\xi=s^{1+a\alpha}\). Then
\(h+P=h+\xi u(Y^a+Z^b)\), while the constraint is
\(s^k(YZ-R)=0\). We write
\(\chi=s^k\chi_{\rm orig}\) for the multiplier of the rescaled
constraint \(YZ-R=0\); the two symbols will be kept distinct.
The rescaled spatial critical equations are
\begin{equation}
\begin{aligned}
h_p&=-\chi R_p,&
\xi(Y^a+Z^b)&=-\chi R_u,\\
\xi a uY^{a-1}&=\chi Z,&
\xi b uZ^{b-1}&=\chi Y.
\end{aligned}
\label{eq:tube-10}
\end{equation}

We first dispose of points over \(p=0\), then of points with
\(y=z=w=0\) and \(p\ne0\). Any remaining hypothetical critical
sequence as \(s\to0\) has, after a subsequence, either
\(p\to p_0\ne0\) or \(p\to0\). We exclude the first alternative
by the cone and root-graph calculations and the second by valuations.

At \(p=0\), \(R_s(0,w)=w^k\) exactly, so the vertical constrained
critical equations are those of \(P\) on \(yz=w^k\).
If \(w\ne0\), then \(y,z\ne0\). Multiplying the \(y\)- and
\(z\)-equations by \(y\) and \(z\), respectively, and eliminating
\(\chi_{\rm orig}w^{k-1}\) from the \(w\)-equation gives
\[
 (ka+1)y^a+z^b=0,\qquad y^a+(kb+1)z^b=0.
\]
Their determinant is \(k(kab+a+b)\ne0\), which would force
\(y^a=z^b=0\), a contradiction. If \(w=0\), then \(yz=0\);
at a nonzero point on either axis \(P_w=y^a+z^b\ne0\), whereas
the \(w\)-derivative of \(yz-w^k\) is zero. Thus the vertical
equations exclude every point over \(p=0\) except the origin,
without a condition on \(w/s\).
At \(y=z=w=0\) with \(p\ne0\), constrained criticality is exactly
criticality of \(d\) on its smooth root graph and is excluded
near the origin.

First exclude critical points whose base coordinate tends to
\(p_0\ne0\) while \(s\to0\). Let
\(\varrho=\max(|y|^{1/\alpha},|z|^{1/\beta},|w|)\).
If \(\varrho/s\to\infty\), weighted rescaling by \(\varrho\) leaves a
nonvertex cone point: every lower term in \(R_s/\varrho^k\) tends to
zero. The normalized vertical critical equations would make
\(P\) critical on the smooth nonvertex cone, which is impossible.
If \(\varrho/s\) is bounded, use the rescaled variables and
multiplier in \eqref{eq:tube-10}. The variables \(Y,Z,u\) are
bounded and \(\xi\to0\). On the bounded \((p,u)\)-set,
\(R_p\) is bounded, while \(h_p(p)\) stays bounded away from zero
because \(h_p(p_0)\ne0\). The base equation \(h_p=-\chi R_p\)
therefore bounds \(\lvert\chi\rvert\) away from zero.
The \(Y\)- and \(Z\)-equations in \eqref{eq:tube-10}, with bounded
\(u,Y,Z\) and \(\xi\to0\), now force \(Y,Z\to0\).
After passing to a subsequence \(u\to u_0\), the constraint gives
\(R(p_0,u_0)=0\). All roots over the selected base cut lie in
the chosen local graph representative; only at this point do we
use its unit component of \(R_p\). The base equation then bounds
\(\chi\) above, so a further subsequence has a nonzero finite
limit \(\chi_{\rm lim}\). The \(u\)-equation in \eqref{eq:tube-10} gives
\(R_u(p_0,u_0)=0\). Together with \(R(p_0,u_0)=0\) and
\(h_p(p_0)=-\chi_{\rm lim}R_p(p_0,u_0)\), this says \(h\) is
critical on the smooth root graph \(R=0\) away from its origin,
contrary to isolation of \(d\).

Thus a remaining critical sequence has \(p\to0\).
Real-analytic curve selection gives an arc; write
\(\nu\) for its valuation, and set
\[
m_0=\nu(p)>0,\qquad
\nu_{\rm v}=\min\bigl(\nu(y)/\alpha,\nu(z)/\beta,\nu(w)\bigr).
\]
The case \(\nu_{\rm v}=\infty\) was already covered by the root graph,
so assume it is finite. Write \(\varrho\) for an arc scale with
\(\nu(\varrho)=\nu_{\rm v}\).
If the lower Weierstrass terms all vanish under rescaling by
\(\varrho\), the vertical equation limits to \(yz=w^k\) at a
nonvertex point. Normalize the original critical multiplier to
\(\chi_{\rm cone}=\chi_{\rm orig}/\varrho^{d_P-k}\), where \(d_P=1+a\alpha\).
A nonzero component of the cone gradient
bounds this multiplier through the corresponding vertical
critical equation; a convergent subsequence of all the equations
then gives \(d(P|_{yz=w^k})=0\), contradicting the
submersivity just proved. Thus, for some \(j<k\),
\begin{equation}
(k-j)(\nu s-\nu_{\rm v})+\nu A_j(p)\le0.
\label{eq:tube-8}
\end{equation}
Since \(A_j(0)=0\), \(\nu A_j(p)\ge m_0\), and therefore
\begin{equation}
\nu_{\rm v}-\nu s\ge m_0/k.
\label{eq:tube-9}
\end{equation}
Return to the rescaled variables and multiplier in \eqref{eq:tube-10}.
By \eqref{eq:tube-9}, \(p,u,Y,Z\to0\) and
\(\nu Y\ge\alpha m_0/k\), \(\nu Z\ge\beta m_0/k\).
The unit component of \(R_p\) and the gradient bound give
\(\nu\chi\le L_0m_0=4m_0\).
If \(Y,Z,u\ne0\), the \(Y\)- and \(Z\)-equations in
\eqref{eq:tube-10} give
\begin{align*}
\nu Z-\nu Y
 &=\nu\xi+\nu u+(a-2)\nu Y-\nu\chi\\
 &\ge \nu\xi+\nu u+
 \left(\frac{(a-2)\alpha}{k}-L_0\right)m_0>0,\\
\nu Y-\nu Z
 &=\nu\xi+\nu u+(b-2)\nu Z-\nu\chi\\
 &\ge \nu\xi+\nu u+
 \left(\frac{(b-2)\beta}{k}-L_0\right)m_0>0.
\end{align*}
Here \(\nu\xi,\nu u>0\), and the strict signs use the two
inequalities on \(a,b\) in Theorem~\ref{thm:tube}, together with
\(m_0>0\). The two conclusions are incompatible.
If exactly one of \(Y,Z\) vanishes, or if \(u=0\) while both
are nonzero, \eqref{eq:tube-10} forces \(\chi=0\), hence \(h_p=0\), contrary
to \(p\ne0\). If both vanish, then \(R=0\) and the remaining
critical equations make \(h\) critical on the root graph,
contrary to isolation of \(d\). This exhausts the arc.
\end{proof}

The outer boundary, the compatible model changes, and isolation are
now established. We next compute the inner collapse and prove acyclicity
of the large total; the local-page comparison will then follow without
any additional geometric assumption on the custom cuts.
